%% characterization -- existential entailment as (I0) and (U) and (II), with
%% the satisfiability dispatch; the arrow example; remarks on covariant and
%% joint-parity instances and on the one-sign case {bot,f}; pure forms of (U)
%% and (II). Closure under intersection (prop:ch-meet) is in classes.tex.

\section{The Characterization}\label{sec:ch}

This section characterizes existential entailment over arbitrary trees by
three conditions on single seeds. Let $X$ and $Z$ be disjoint finite sets of
variables, let $\varphi$ be a system over $X$, and let $\psi$ be a system
over $X\cup Z$ (\cref{def:bg-existential}); every variable of $\psi$ outside
$Z$ belongs to $X$. An atom $(v^s,w)$ is \emph{visible} if $v\in X$ and
\emph{hidden} if $v\in Z$. Restricting an assignment to $X$ restricts its set
of atoms to $\Conf_X$, so by \cref{lem:lf-clauses}
\begin{equation}\label{eq:ch-models}
 \varphi\models\Exists\psi\quad\Longleftrightarrow\quad
 \Mod(\varphi)\subseteq\mathcal F_\psi,\qquad\text{where }
 \mathcal F_\psi=\{T|_X:T\in\Mod(\psi)\}.
\end{equation}
Here $\Mod(\varphi)$ consists of sets of visible atoms, and $\mathcal F_\psi$
is the \emph{projected family} of $\psi$.

\subsection{Three Conditions}\label{sec:ch-theorem}

\begin{definition}[The three conditions]\label{def:ch-conditions}
For $\varphi$ and $\psi$ as above, consider:
\begin{description}[leftmargin=2.8em,font=\normalfont]
\item[\CondI] $\mu_\psi|_X\subseteq\mu_\varphi$;
\item[\CondU] for every visible atom $a$ that is live in $\varphi$,
  $\Reach_\psi(a)$ contains no forbidden root of $\psi$;
\item[\CondII] for every visible atom $a$ that is live in $\varphi$,
  $\Reach_\psi(a)|_X\subseteq\Cl_\varphi(a)$.
\end{description}
\end{definition}

The names are mnemonic: \CondI\ is the inclusion for zero seeds, comparing the
consequences of the facts alone; \CondU\ concerns the unit clauses that are
forbidden roots; and \CondII\ is the inclusion for a single seed $a$.

When both systems are satisfiable, $\mu_\psi$ and $\mu_\varphi$ are their
least models and $\Cl_\varphi(a)$ is the least model of $\varphi$ that
contains $a$ (\cref{lem:lf-closure}(2)). So \CondI\ says that every visible
atom forced by $\psi$ is forced by $\varphi$. Condition \CondU\ says that
$\psi$ never derives a forbidden root from an atom that some solution of
$\varphi$ contains. Condition \CondII\ says that every solution of $\varphi$
that contains an atom $a$ also contains every visible atom that $\psi$
derives from $a$.

\begin{theorem}[Characterization]\label{thm:ch-characterization}
Let $\varphi$ be a system over $X$ and $\psi$ a system over $X\cup Z$, for
any arity and any variance, and let assignments range over arbitrary trees.
\begin{enumerate}[leftmargin=*]
\item If $\varphi$ is unsatisfiable, then $\varphi\models\Exists\psi$. If
  $\varphi$ is satisfiable and $\psi$ is not, then
  $\varphi\not\models\Exists\psi$.
\item If $\varphi$ and $\psi$ are satisfiable, then
  $\varphi\models\Exists\psi$ iff \CondI, \CondU\ and \CondII\ hold.
\end{enumerate}
Equivalently, $\Eall(\varphi,\psi)\iff\neg\Sat(\varphi)\lor
\bigl(\Sat(\psi)\land\CondI\land\CondU\land\CondII\bigr)$.
\end{theorem}

Each condition concerns a single seed. An obstruction to extending a model
$S$ of $\varphi$ may a priori involve two atoms of $S$ whose closures under
$\psi$ collide in an exclusion pair, possibly at a hidden variable. The
proof below uses collision return (\cref{lem:lf-collision}) to turn every
such collision into a violation of \CondI\ or of \CondII\ at one seed. This
step uses duality, which holds for every variance; no condition on pairs of
seeds remains. When both systems are satisfiable, \cref{fig:ch-conditions}
shows the four kinds of obstruction that the proof meets: three of them are the failures of the three
conditions, and the fourth, a collision of two seeds, is reduced to the
first or the third.

\begin{figure}[tbp]
\centering
\begin{tikzpicture}[>=Stealth,font=\small,x=1cm,y=1cm,
    every node/.style={inner sep=1.5pt},
    model/.style={draw,rounded corners=3pt,fill=black!8},
    pt/.style={circle,fill,inner sep=0pt,minimum size=3.5pt},
    der/.style={->,decorate,decoration={snake,amplitude=0.6pt,segment length=5pt,post length=3pt}},
    dual/.style={->,dashed},
    opp/.style={densely dotted,thick},
    lb/.style={font=\scriptsize,inner sep=1pt},
    ttl/.style={font=\footnotesize,anchor=south}]
  \foreach \k in {0,1,2,3} {
    \begin{scope}[xshift=3.15*\k cm]
      \draw[black!50] (0,0) -- (2.95,0);
      \draw[black!25] (0,-1.25) rectangle (2.95,1.6);
    \end{scope}}
  \node[lb,anchor=east,align=right] at (-0.05,0.8) {visible\\ atoms};
  \node[lb,anchor=east,align=right] at (-0.05,-0.6) {hidden\\ atoms};
  % (a) (I0)
  \node[ttl] at (1.475,1.6) {(a) \CondI\ fails};
  \draw[model] (0.2,0.2) rectangle (1.6,1.4);
  \node[lb] at (0.55,1.2) {$\mu_\varphi$};
  \node[pt,label={[lb]below:fact $e$}] (ea) at (0.9,-0.7) {};
  \node[pt,label={[lb]right:$b$}] (ba) at (2.4,0.75) {};
  \draw[der] (ea) -- (ba);
  % (b) (U)
  \begin{scope}[xshift=3.15cm]
    \node[ttl] at (1.475,1.6) {(b) \CondU\ fails};
    \draw[model] (0.2,0.2) rectangle (1.6,1.4);
    \node[lb] at (0.6,1.2) {$\Cl_\varphi(a)$};
    \node[pt,label={[lb]left:$a$}] (ab) at (0.9,0.65) {};
    \node[draw,double,rounded corners=1pt,lb,inner sep=2pt,align=center] (fb) at (2.2,-0.7) {forbidden\\ root};
    \draw[der] (ab) -- (fb);
  \end{scope}
  % (c) (II)
  \begin{scope}[xshift=6.3cm]
    \node[ttl] at (1.475,1.6) {(c) \CondII\ fails};
    \draw[model] (0.2,0.2) rectangle (1.6,1.4);
    \node[lb] at (0.6,1.2) {$\Cl_\varphi(a)$};
    \node[pt,label={[lb]left:$a$}] (ac) at (0.9,0.65) {};
    \node[pt,label={[lb]right:$b$}] (bc) at (2.4,0.75) {};
    \draw[der] (ac) .. controls (1.6,-0.9) and (2.2,-0.6) .. (bc);
  \end{scope}
  % (d) collision
  \begin{scope}[xshift=9.45cm]
    \node[ttl] at (1.475,1.6) {(d) two seeds collide};
    \draw[model] (0.15,0.2) rectangle (1.75,1.4);
    \node[lb] at (0.45,1.2) {$S$};
    \node[pt,label={[lb]above:$a$}] (ad) at (0.5,0.6) {};
    \node[pt,label={[lb]above:$(q,v)$}] (qd) at (1.35,0.85) {};
    \node[pt,label={[lb]below:$b$}] (bd) at (1.0,-0.65) {};
    \node[pt,label={[lb]below:$\bar b$}] (nbd) at (2.0,-0.65) {};
    \node[pt,label={[lb]above:$(\bar q,v\delta)$}] (rd) at (2.4,0.95) {};
    \draw[der] (ad) -- (bd);
    \draw[der] (qd) -- (nbd);
    \draw[opp] (bd) -- (nbd);
    \draw[dual] (bd) .. controls (1.5,0.0) and (2.1,0.4) .. (rd);
  \end{scope}
\end{tikzpicture}
\caption{When both systems are satisfiable, the four ways in which a model $S$
of $\varphi$ can fail to extend to a model of $\psi$, in the proof of
\cref{thm:ch-characterization}. Upper halves
hold visible atoms, lower halves hidden ones; wavy arrows are derivations of
$\psi$, and the shaded box is a model of $\varphi$. (a) A fact of $\psi$
derives a visible atom outside the least model $\mu_\varphi$. (b) A seed $a$
derives a forbidden root; then the least model $\Cl_\varphi(a)$ of $\varphi$
that contains $a$ cannot be extended. (c) A seed $a$ derives a visible atom
outside $\Cl_\varphi(a)$. (d) Two seeds of $S$ derive opposite atoms $b$ and
$\bar b$, possibly hidden. Collision return (\cref{lem:lf-collision}) reverses
the second derivation by duality (dashed): the single seed $a$ derives the
opposite of the extension $(q,v\delta)\in S$ of the second seed, an atom
outside $S\supseteq\Cl_\varphi(a)$, so this is case~(c). If $a$ is a fact of
$\psi$ instead, it is case~(a).}
\Description{Four panels, each split by a horizontal line into visible atoms
above and hidden atoms below. In each panel a shaded box is a model of phi.
Panel a: a fact e of psi, below the line, derives by a wavy arrow a visible
atom b outside the box labeled mu phi. Panel b: a seed a inside the box,
labeled closure of a, derives a doubly framed forbidden root below the line.
Panel c: a seed a inside the box derives, through hidden atoms, a visible atom
b outside the box. Panel d: two seeds a and q v inside the box labeled S
derive the hidden atoms b and b bar, joined by a dotted line; a dashed arrow from b
leads to the visible atom q bar, v delta outside the box.}
\label{fig:ch-conditions}
\end{figure}


\begin{proof}
Item (1) is immediate from \cref{def:bg-existential}.

For (2), assume that $\varphi$ and $\psi$ are satisfiable. We first show
that the three conditions are necessary. Assume $\varphi\models\Exists\psi$.
The least model $\mu_\varphi$ of $\varphi$ is $T|_X$ for some model $T$ of
$\psi$ by \cref{eq:ch-models}, and $T\supseteq\mu_\psi$ by
\cref{lem:lf-closure}(2); hence $\mu_\psi|_X\subseteq T|_X=\mu_\varphi$,
which is \CondI. Let $a$ be a visible atom that is live in $\varphi$. Its
least model $\Cl_\varphi(a)$ is $T|_X$ for a model $T$ of $\psi$. Since $T$
contains $a$ and is closed under the rules of $\psi$, it contains $\Reach_\psi(a)$.
As $T$ contains no forbidden root, this gives \CondU, and
$\Reach_\psi(a)|_X\subseteq T|_X=\Cl_\varphi(a)$ gives \CondII.

Conversely, assume \CondI, \CondU\ and \CondII, let $S$ be a model of
$\varphi$, and put $T=H_\psi(S)$, where $S$ is read as a set of atoms over
$X\cup Z$. By \cref{lem:lf-closure}(5) and \cref{eq:ch-models}, it suffices
to show that $T|_X=S$ and that $T$ is consistent. Every atom of $S$ is live
in $\varphi$, and $\Cl_\varphi(a)\subseteq S$ for every $a\in S$ by
\cref{lem:lf-closure}(1), since $S$ is closed under the rules of $\varphi$ and
contains its facts.

\emph{Visible part.} Clearly $S\subseteq T|_X$. A visible atom of $T$ lies in
$\mu_\psi|_X\subseteq\mu_\varphi\subseteq S$, by \CondI, or in $\Reach_\psi(a)|_X$
for some $a\in S$, which lies in $\Cl_\varphi(a)\subseteq S$ by \CondII.

\emph{Consistency.} Suppose that $T$ is not consistent, and take seeds
$x,y\in\Fact_\psi\cup S$ as in \cref{lem:lf-closure}(3). If $\Reach_\psi(x)$
contains a forbidden root, then $x\notin\Fact_\psi$, because $\mu_\psi$ is
consistent; so $x\in S$, contradicting \CondU. Otherwise $b\in \Reach_\psi(x)$
and $\bar b\in \Reach_\psi(y)$ for some atom $b$; since an exclusion pair is
unordered, we may assume that $y\in S$ if one of $x,y$ lies in $S$.
\begin{itemize}[leftmargin=*]
\item If $x,y\in\Fact_\psi$, then $b,\bar b\in\mu_\psi$, contradicting the
  consistency of $\mu_\psi$.
\item If $x\in\Fact_\psi$ and $y=(q,v)\in S$, then \cref{lem:lf-collision}
  gives $(\bar q,v\delta)\in \Reach_\psi(x)\subseteq\mu_\psi$ for some $\delta$.
  This atom is visible, so it lies in $\mu_\varphi\subseteq S$ by \CondI. But
  $(q,v\delta)\in S$ because $S$ is extension closed, and $S$ has an
  exclusion pair.
\item If $x\in S$ and $y=(q,v)\in S$, possibly $x=y$, then
  \cref{lem:lf-collision} gives $(\bar q,v\delta)\in \Reach_\psi(x)|_X$ for some
  $\delta$, which lies in $\Cl_\varphi(x)\subseteq S$ by \CondII. Again
  $(q,v\delta)\in S$, and $S$ has an exclusion pair.
\end{itemize}
Each case contradicts an assumption. Hence $T$ is a model of $\psi$ with
$T|_X=S$, so $S\in\mathcal F_\psi$.
\end{proof}

The proof shows that, when the conditions hold, $H_\psi(S)$ is the least
extension of every model $S$ of $\varphi$; when a condition fails, the least
model $\mu_\varphi$ or $\Cl_\varphi(a)$ that refutes it is a countermodel.
The two qualifications of the theorem are needed.
\begin{itemize}[leftmargin=*]
\item \emph{The satisfiability dispatch.} Let $X=\{x\}$, $\varphi$ empty, and
  $\psi=\{z=\Bot,\ \Top\le z\}$ with $z$ hidden. Then $\mu_\psi$ consists of
  atoms of $z$, and $\psi$ has no rule that involves $x$ other than its
  extensions, so the three conditions hold; but $\psi$ is unsatisfiable, and
  the entailment fails.
\item \emph{Condition \CondU.} Let $n\ge0$, $X=\{y\}$, $\varphi$ empty, and
  $\psi=\{y\le f(z_i)_i\}$ with distinct hidden variables $z_i$. Then $\psi$
  has no facts, and its only pure moves from visible atoms are the
  pops $(y^0,iw)\to(z_i^{c(i)},w)$, after which no pure move applies. By
  \cref{lem:ch-pure} below, \CondI\ and \CondII\ hold. But $y=\Top$ has no
  extension: \CondU\ fails at the live atom $(y^0,\eps)$, which is a
  forbidden root of $\psi$.
\end{itemize}
The theorem also depends on duality. For a set of one-premise rules without
duals, say $(x^0,w)\to(z^0,w)$ and $(y^0,w)\to(z^1,w)$ with $z$ hidden, the
three conditions hold for the empty antecedent, but $x=y=\Top$ produces an
exclusion pair at $z$. No system has exactly these rules: by
\cref{lem:lf-duality}, a system with these rules also has the duals
$(z^1,w)\to(x^1,w)$ and $(z^0,w)\to(y^1,w)$, which expose the collision as the
visible consequences $(x^0,w)\Rightarrow^*(y^1,w)$ and
$(y^0,w)\Rightarrow^*(x^1,w)$, and \CondII\ fails.

The conditions only need pure derivations of $\psi$, from extended seeds.

\begin{lemma}[Pure forms]\label{lem:ch-pure}
\leavevmode
\begin{enumerate}[leftmargin=*]
\item \CondU\ holds iff there are no visible atom $(p,u)$ that is live in
  $\varphi$, word $\alpha$ and forbidden root $e$ of $\psi$ with
  $(p,u\alpha)\to^*_\psi e$.
\item \CondII\ holds iff $b\in\Cl_\varphi(p,u)$ for every visible atom
  $(p,u)$ that is live in $\varphi$ and every visible atom $b$ with
  $(p,u)\to^*_\psi b$.
\end{enumerate}
\end{lemma}

\begin{proof}
If $(p,u)$ is live in $\varphi$, so is $(p,u\alpha)$, and
$\Cl_\varphi(p,u\alpha)\subseteq\Cl_\varphi(p,u)$
(\cref{lem:lf-closure}(4)). By \cref{lem:lf-ext-first}, $b\in \Reach_\psi(p,u)$
iff $(p,u\alpha)\to^*_\psi b$ for some $\alpha$. Item (1) follows directly,
since $(p,u\alpha)$ is again a live visible atom. For (2), the condition is
clearly necessary. Conversely, if $b\in \Reach_\psi(p,u)|_X$ for a live $(p,u)$,
then $(p,u\alpha)\to^*_\psi b$ for some $\alpha$, and the condition, applied
to the live atom $(p,u\alpha)$, gives
$b\in\Cl_\varphi(p,u\alpha)\subseteq\Cl_\varphi(p,u)$.
\end{proof}

\begin{example}[The consequent $x\le z\to z$]\label{ex:ch-arrow}
For the arrow of \cref{ex:lf-arrow}, let $X=\{x\}$, $Z=\{z\}$ and
$\psi=\{x\le z\to z\}$, the literal $x\le f(z,z)$. By
\cref{tab:lf-clauses}, $\psi$ has the forbidden root $(x^0,\eps)$, no
facts, and the pure moves
\[
 (x^0,0w)\to(z^1,w)\to(x^1,1w),\qquad (x^0,1w)\to(z^0,w)\to(x^1,0w).
\]
No pure move leaves a negative atom of $x$. So the pure derivations of $\psi$
of positive length between visible atoms are exactly
$(x^0,iw)\to^*(x^1,\bar\imath w)$, where $\bar\imath=1-i$: $\psi$ exchanges
the first letter and the sign. Since $\mu_\psi=\varnothing$, \CondI\ holds for
every antecedent, and $\psi$ is satisfiable, with $x=\Bot$.

\emph{The empty antecedent.} Let $\varphi_0$ be empty. Then \CondU\ fails at
$(x^0,\eps)$, with the countermodel $x=\Top$. Condition \CondII\ fails at
$(x^0,0)$: $\psi$ derives $(x^1,1)$, but $\Cl_{\varphi_0}((x^0,0))$ consists of
the atoms $(x^0,0w)$, the leaf pair $(0\Sigma^*,\varnothing)$ of the
countermodel $x=\Bot\to\Omega$ (\cref{ex:lf-arrow}). Indeed
$\Bot\to\Omega\le_c z\to z$ would require $z\le_c\Bot$ and $\Omega\le_c z$.
The tree $\Bot\to\Top$, with leaf pair $(\Sigma^+,\varnothing)$, fails for a
reason that a priori involves two seeds: the closures of $(x^0,0)$ and
$(x^0,1)$ under $\psi$ contain the exclusion pair $(z^1,\eps)$ and
$(z^0,\eps)$. Collision return reverses the second branch
$(x^0,1)\to(z^0,\eps)$ into $(z^1,\eps)\to(x^1,1)$, which continues the first
branch: the single seed $(x^0,0)$ derives $(x^1,1)$, the opposite of the
second seed, and \CondII\ fails at $(x^0,0)$. The extension in
\cref{lem:lf-collision} is needed in general. The atom $(x^0,\eps)$ derives
$(z^0,\eps)$ through its extension $(x^0,1)$, which collides with the atom
$(z^1,\eps)$ derived from $(x^0,0)$; collision return gives $(x^1,1)$, the
opposite of an extension of $(x^0,\eps)$, whereas $(x^1,\eps)$ is not
derivable from $(x^0,0)$.

\emph{The antecedent $x\le x\to x$.} Let $\varphi_1=\{x\le x\to x\}$ be the
system of \cref{ex:lf-clauses}. It is satisfiable. Condition \CondU\ holds,
because $(x^0,\eps)$ is a forbidden root of $\varphi_1$ as well, hence dead.
For \CondII, by \cref{lem:ch-pure} it suffices that
$(x^1,\bar\imath w)\in\Cl_{\varphi_1}((x^0,iw))$ for every live $(x^0,iw)$,
and $\varphi_1$ derives it: $(x^0,0w)\to(x^1,w)\to(x^1,1w)$ by $\mathrm{pop}_0$ and
$\mathrm{push}_1$, and $(x^0,1w)\to(x^0,w)\to(x^1,0w)$ by $\mathrm{pop}_1$
and $\mathrm{push}_0$. By \cref{thm:ch-characterization},
$\varphi_1\models\exists z\,\psi$; indeed $z=x$ is a witness.
\end{example}

\subsection{Special Cases}\label{sec:ch-remarks}

\begin{remark}[Covariant and joint-parity instances]\label{rem:ch-parity}
A \emph{parity assignment} of a system $\chi$ is a map $\pi:V\to\{0,1\}$ with
$\pi(u)=\pi(v)$ for every literal $u\le v$ and $\pi(a_i)=\pi(u)\oplus c(i)$
for every $i$ and every literal $f(a_i)_i\le u$ or $u\le f(a_i)_i$ of $\chi$.
By \cref{tab:lf-clauses}, every rule of $\chi$ then preserves the bit
$\kappa(v^s,w)=s\oplus\pi(v)$ of an atom; for instance, $\mathrm{push}_i$ of
$f(a_i)_i\le u$ leads from $\kappa=c(i)\oplus\pi(a_i)=\pi(u)$ to
$\kappa=\pi(u)$. Opposite atoms have different bits. Hence no derivation
leads from an atom to the opposite of one of its extensions: self-collisions,
cause (c) of \cref{cor:lf-dead}, do not occur, and in a satisfiable system
every dead atom is dead by support. Likewise, in \cref{lem:lf-collision}
the two colliding seeds have different bits, so in the proof of
\cref{thm:ch-characterization} two branches from the same seed never collide.

More is true. Rename every atom $(v^s,w)$ to $(v^{s\oplus\pi(v)},w)$. By
\cref{tab:lf-clauses}, this maps the rules, facts and forbidden
roots of $\chi$ onto those of the following system $\chi^\pi$ for the
covariant variance $c=0$: $\chi^\pi$ keeps every literal of $\chi$ whose root
has $\pi=0$ and reverses every other one, where the root of $u\le v$, $u=\Bot$
and $\Top\le u$ is $u$. Reversal turns $u\le v$ into $v\le u$, a lower
constructor literal into the upper one with the same variables and
conversely, and exchanges $u=\Bot$ and $\Top\le u$. On trees the renaming is
$\rho(v)\mapsto\nu_{\pi(v)}\rho(v)$, read covariantly: the covariant leaf
cone of sign $s\oplus\pi(v)$ of $\nu_{\pi(v)}\rho(v)$ is the cone $L^s$ of
$\rho(v)$, because $(\nu_rt)^\partial=\nu_{1\oplus r}t$ by
\cref{lem:bg-normalization}(1). So $\rho$ solves
$\chi$ iff this assignment solves $\chi^\pi$, by \cref{lem:lf-clauses}. If
$\varphi\wedge\psi$ has a \emph{joint} parity assignment, one renaming serves
both systems and keeps visibility, and the instance is equivalent to a
covariant instance of the same size, in every tree class, since $\nu_r$ is an
involution that preserves finiteness and regularity
(\cref{lem:bg-normalization}). The covariant case itself is $\pi=0$, where
every rule keeps the sign and the rules of sign $1$ are the duals of the
rules of sign $0$.

\Cref{thm:ch-characterization} needs no parity assignment, and the system
$x\le x\to x$ has none, since it would require $\pi(x)=\pi(x)\oplus c(0)$. Its
self-collision $(x^0,00)\to(x^1,0)$ (\cref{ex:lf-least}) is exactly what a
parity assignment excludes.
\end{remark}

\begin{remark}[One sign: the signature $\BF$]\label{rem:ch-one-sign}
Let $f$ be covariant, and consider trees over $\BF$ and systems without
literals $\Top\le u$ (\cref{sec:bg-trees}). Every leaf of such a tree has
sign $1$, so $L^0_t=\varnothing$, and $L^1_t$ is the complement of the set
$\Kpos(t)$ of positions labeled $f$ (\cref{thm:lf-representation}(4)). The
trees over $\BF$ thus correspond to the extension-closed sets of negative
atoms. All facts of such a system come from literals $u=\Bot$ and are
negative atoms, and every rule keeps the sign, so the closure of a set of negative atoms consists of
negative atoms, and an exclusion pair never arises. The models over $\BF$ are
therefore the sets of negative atoms that contain the facts, are
closed under the rules, and avoid the forbidden roots $(u^1,\eps)$ of the
lower constructor literals $f(a_i)_i\le u$. \Cref{lem:lf-closure} and
\cref{thm:ch-characterization} hold over $\BF$, with satisfiability over
$\BF$, by the same proofs, in which the exclusion case is void, once every
set of atoms is taken in the universe
$\Conf^-_V=\{v^1:v\in V\}\times\Sigma^*$ of negative atoms: the seed sets,
the visible sets, the models, the closures, liveness, and the atoms
quantified in the three conditions are all negative. (Without this
restriction the statement fails: for the empty system the closure of
$(x^0,\eps)$ is consistent, but no $\Top$-free assignment contains it.)

Read through complements, a negative atom $(v^1,w)$ says $w\notin\Kpos(v)$,
and each rule on negative atoms is the contrapositive of a clause on
constructor positions: $u\le v$ says $\Kpos(u)\subseteq\Kpos(v)$,
$f(a_i)_i\le u$ says $\eps\in\Kpos(u)$ and $i\Kpos(a_i)\subseteq\Kpos(u)$,
$u\le f(a_i)_i$ says that $iw\in\Kpos(u)$ implies $w\in\Kpos(a_i)$, $u=\Bot$
says $\eps\notin\Kpos(u)$, and extensions say that $\Kpos(v)$ is prefix closed.
In these terms, \CondI\ says that every visible position that some solution
of $\varphi$ puts into a variable is put there by some solution of $\psi$;
\CondU\ says that every visible position that every solution of $\psi$ puts
into a variable is put there by every solution of $\varphi$; and \CondII,
read contrapositively, says that every solution of $\varphi$ is closed under
the visible consequences of the clauses of $\psi$. Over $\BF$ the three
conditions therefore compare, respectively, the unions of all solutions, the
least solutions, and the closures of single positions, for $\varphi$ and for
the projection of $\psi$. In this picture the solution sets are closed under
nonempty unions and intersections of constructor positions, no exclusions
and no self-collisions occur, and the only obstructions are single seeds. The
decision procedure of this paper does not treat $\BF$ separately; it reduces
$\BF$ to $\BFT$ (\cref{sec:tc}).
\end{remark}

A forbidden root has the empty word, so the pure derivations in
\cref{lem:ch-pure}(1) end with an empty stack. A failure of \CondII\ with a
live seed $a$ and a consequence $b$ is thus witnessed by a pure derivation of
$\psi$ from a live extension of $a$ to $b$, with $b$ outside the least model
of $\varphi$ that contains that extension. \Cref{sec:sm,sec:un,sec:ex} turn the satisfiability tests, \CondI,
\CondU\ and this form of \CondII\ into problems about finite automata.
